\cotefile{19}
\scanpage{65}
\mspage{5 ter}
% archive sheet number: 64

\(\unsure{,}\;\unsure{6|b}\;=\)

% slanted
\margin{\uline{Below}}

\noindent\llap{\(\left[\vphantom{\rule{0pt}{3em}}\right.\)\quad}\uline{Assume} moreover that \(i_{\unsure{+|*}}\) exists equally. Then \unsure{Mitchell} \unsure{Imbedding|structure|situation} theorem applies. Moreover, \uuline{in this case} equivalent conditions

\diagram{\begin{tikzcd}[ampersand replacement=\&,column sep=small,row sep=small]
\text{a) }j^{!}_{\deleted{\unsure{!}}}\ \text{exact}
  \arrow[r,Rightarrow,"{\text{b)}}"] \&
i_*\ \text{exact}\arrow[r,Rightarrow,"{\text{c)}}"]\arrow[d,Rightarrow] \&
\varphi=j^{*}i_*\ \text{is exact}\arrow[d,Rightarrow,bend left=22] \\
\& j_*\ \text{\unsure{resp|Ker} \unsure{épaisses|épaisseurs|épaisse}}\arrow[d,Rightarrow] \& {} \\
\& F\longrightarrow i_{\unsure{*|+}}{\unsure{j|i}}^{*}F\ \text{surjectif} \&
\end{tikzcd}}
\scanpage{66}
% archive sheet number: 65

\deleted{\uline{\unsure{Question}} Is it enough that \(\unsure{j|j^*}\) \illegible{2} Or that \(i_!\) be exact?}

\deleted{\(i_!i^*F\longrightarrow F\) has a kernel \(N\) such that \(i^*N=0\)\unsure{?} i.e. of the type \(j_*H\), \deleted{\unsure{i.e.}} take \(H\unsure{=|\simeq}j^*N\), and \(j^*\) being exact, we have \deleted{\(\longrightarrow N\longrightarrow i_!i^*F\longrightarrow F\), or \(j^*N\longrightarrow\)} \(\longrightarrow j^*N\longrightarrow j^*i_!i^*F\longrightarrow j^*F\). Yet \(j^*i_!=0\), hence \(j^*N=0\) hence \(N=0\). \(\,i_!i^*N\longrightarrow N\)}

\uline{Question}\unsure{.|:} Is it enough \(i_!\) exist ?? \uuline{No}, as we see in the case dual to:

\diagram{\begin{tikzcd}[ampersand replacement=\&,column sep=large,row sep=tiny]
{\substack{C''=A\text{-}\unsure{\text{mods}|\text{mds}|\text{modules}}\\\text{killed by some}\\ f\in S}}
\arrow[r,shift left,"{j_*}"] \&
C=C_A \arrow[l,shift left,"{\substack{j_!\\\text{biggest}\\\text{submodule}\\\text{locally}\\\text{killed}\\\text{by elem.}\\\text{of }S}}"]
\arrow[d]\arrow[r,shift left,"{i^*\ \text{localization}}"] \&
C'=C_{S^{-1}A}\arrow[l,shift left,"{i_*\ \substack{\text{restriction}\\\text{of scalars}}}"']\\
\unsure{1} \& A\ \text{comm. ring} \&
\end{tikzcd}}

\(i_*\) is exact, but not \(j_!\) in general (\(A\) \inserted{discrete} valuation ring, \(S=\{\pi\}\). \(\pi\) uniformizing)
\scanpage{67}
\mspage{6}
% archive sheet number: 66

\(\ast\) foncteur (exact, surjectif i.e. un passage au quotient) \inserted{\(i^*\)} qui a \uline{deux} adjoints \(i_!,i_*\) (dont l'un \inserted{\(i_!\)} \uline{exact}, \unsure{\}} et \(=\) \unsure{plein fid.|pleinement fid.}: \uline{\unsure{im épaisse|une ss épaisse|image épaisse}})

% slanted
\margin{\unsure{dual}}

\(\ast\) foncteur (exact, \unsure{pleinemt|pleinement} fid. et \(=\): image épaisse) \(j_*\) qui a deux adjoints \inserted{\(j^*,j^!\)} (dont l'un \inserted{\(j^*\)} est exact, \deleted{surjectif}, \unsure{ie|et} \(=\) un passage au quotient)

\(\ast\) Situations mixtes \(i^*_{\unsure{!,*|*}},j^{\unsure{a|*}},j_*\) et autres éventuellement

\(\ast\) Foncteur exact à gauche

\noindent\rule{4em}{0.4pt}

Start with situation

\diagram{\begin{tikzcd}[ampersand replacement=\&,column sep=small]
C\arrow[r,shift left,"{i^*}"] \& C'\arrow[l,shift left,"{i_*}"] \&
\deleted{\unsure{\text{exact}}},\ i_*\ \text{fully faithful}\ \deleted{[\text{\unsure{and image thick}}]}\arrow[d,Rightarrow]\\
\& \& i^*\ \deleted{\unsure{\text{exact}}}\ \text{passage to quotient}\ \deleted{[\text{\unsure{and exact}}]}
\end{tikzcd}}

Then the following conditions are equivalent

\diagram{\begin{tikzcd}[ampersand replacement=\&,row sep=small]
\parbox{.85\linewidth}{a) \(i^*\) \deleted{\unsure{exact}} has another adjoint \(i_!\), \(i_!\) \deleted{\uline{is exact}} \inserted{has thick} image \inserted{\([\unsure{\Leftrightarrow|\mathrm{N.B.}}\ i_!i^*F\longrightarrow F\ \text{\unsure{surj.|syr.} et }i_*\text{ ou }i_!\ \text{pl. fid.}]\)}}\arrow[d,Leftrightarrow]\\
\parbox{.85\linewidth}{b) \(i^*\) is \deleted{\unsure{right}} \uline{exact} [then of the type \(C\longrightarrow C/j_*C'\), where \(j_*\) is the imbedding of a thick subcategory \(C''\) of \(C\)] \deleted{\unsure{and}} \(j^*\) exist and is exact}\arrow[d,Leftrightarrow]\\
\parbox{.85\linewidth}{c) \unsure{Équivalence|equivalent} with the ``standard'' situation defined by a \inserted{left exact} functor \(\varphi\)}
\end{tikzcd}}

% slanted
\margin{Symmetric point of view Start with
\(\begin{tikzcd}[ampersand replacement=\&,column sep=tiny] C''\arrow[r,shift right,"{j_*}"']\& C\arrow[l,shift right,"{j^*}"']\end{tikzcd}\)
s. th. \(j_*\) fully faithful \(\Longleftrightarrow j^*\) passage to quotient. \unsure{équiv.} conditions \(\unsure{j^*|j_*}\) exact, a) \(j_*\) has thick image, and defining \(C'=C/C''\) and \(i^*:C\longrightarrow C'\), \(i^*\) has an adjoint \(i_*\) ???}

We have \deleted{already checked} \(\begin{smallmatrix}c)\Rightarrow a)\\ \phantom{c)}\Rightarrow b)\end{smallmatrix}\) we have checked already even \(b)\Rightarrow c)\), now let us check \(a)\Rightarrow b)\).

Of course if we assume \(a)\), then \(i^*\) is exact (having two adjoints)\unsure{.} Introduce \(C''\) as in \(b)\). We want to define \(j^*F\) as

\[
j^*F=\operatorname{Coker}(i_!i^*F\longrightarrow F)
\]
\scanpage{68}
% archive sheet number: 67
% Writing is inverted relative to the archive foliation. Long oblique strokes cancel the working.

\diagram{\begin{tikzcd}[ampersand replacement=\&,column sep=small]
\deleted{C''}\arrow[r,shift right,"{j_*}"']\&
\deleted{C}\arrow[l,shift right,"{j^*}"']\arrow[r,shift right,"{i^*}"']\&
\deleted{C'}\arrow[l,shift right,"{i_!}"']
\end{tikzcd}}

\[
\begin{aligned}
&\deleted{\unsure{j^*j_*}=\unsure{1_{C''}}}\\
&\deleted{\unsure{i_!i^*}=\unsure{1_{C'}}}
\end{aligned}
\]

\[
\deleted{\left\{\begin{aligned}&\underline{\Gamma}_Y(F)\\&F^U\end{aligned}\right.}
\]

\[
\deleted{\left\{\begin{aligned}
0&\longrightarrow F^U-f-F^Y-0\\
0&\longrightarrow\underline{\Gamma}_Y F\longrightarrow F\unsure{{}^{\wedge}\!{-}|\longrightarrow}\underline{\Gamma}_U \unsure{f|F}
\end{aligned}\right.}
\qquad
\deleted{\begin{aligned} f_!f^* &&f=i,j\\ f_*f^! &&f=i,j\end{aligned}}
\]

\diagram{\begin{tikzcd}[ampersand replacement=\&,column sep=large,row sep=small]
\deleted{C''}\arrow[r,"{j_*}" description] \& \deleted{C}\arrow[l,shift right=2,"{j^*}"']\arrow[l,shift left=2,"{j^!}"]\arrow[r,"{i^*}" description] \& \deleted{C'}\arrow[l,shift right=2,"{i_!}"']\arrow[l,shift left=2,"{i_*}"]\\
\deleted{\unsure{j_*j^*i_*}}\ \deleted{\varphi} \& {}\arrow[l,draw=none]\& {}\arrow[l,shift right,"{j_*\varphi}"]\arrow[l,shift left,"{\deleted{\unsure{j_*\varphi}}}"']\arrow[ll,bend left=12,"{\varphi=j^*i_*}"']
\end{tikzcd}}

\deleted{\unsure{two}} \deleted{\unsure{alltogether|alltoghether|altogether} 12 functors (\unsure{non-trivial} \unsure{ones})}

\diagram{\(\deleted{\unsure{j^*j_*j}}\quad
\left\{\begin{array}{cc}
\deleted{j_*\quad j^*\quad j^!}&\\
\deleted{i^*\quad i_!\quad i_*}&\\ \cline{1-1}
\deleted{\unsure{j}}\quad\deleted{i_!i^*,i_*i^*}&\deleted{C\unsure{\twoheadrightarrow|\rightarrow}C}\\
\deleted{j_*j^*,j_*j^!}&\\ \cline{1-1}
\deleted{\varphi=j^*i_*}&\deleted{C'\longrightarrow C''}\\
\deleted{j_*\varphi}&\deleted{C'\longrightarrow C}
\end{array}\right.\)}
\scanpage{69}
% archive sheet number: 68

\uline{\unsure{Noter|Notes|Notons}} \((i^*,j^*)\) commutent aux \(\underrightarrow{\lim}\) quelconques \inserted{si elles existent dans \(C',C''\)}, et \(=\) pour que \(\underrightarrow{\lim}_{i}\,(G_i,H_i,\unsure{U_i|u_i|U_i:|\lambda_i})\) existe, il suffit que \(\underrightarrow{\lim}_{i}G_i\) et \(\underrightarrow{\lim}_{i}H_i\) existent, et alors \(\ldots\)

Pour les propriétés d'exactitude \uline{à gauche}, \((i^*,j^*)\) est \deleted{\unsure{exactement|extrêmement}} ``aussi bon que'' \(\varphi:C'\longrightarrow C''\), p.ex. si \(\varphi\) commute aux \(\underleftarrow{\lim}\) finies (resp. quelconques) et si \deleted{\unsure{pour}} \inserted{un} type \deleted{une telle} de \(\underleftarrow{\lim}\) existe \unsure{en|dans} \(C',C''\), alors il existe dans \(C\) et \unsure{les|des} foncteurs \((i^*,j^*)\) y commutent \(\ldots\)
\scanpage{71}\mspage{1}
% archive sheet number: 70

\diagram{\begin{tikzcd}[ampersand replacement=\&] A\arrow[r,shift left,"f"]\&B\arrow[l,shift left,"g"]\end{tikzcd}}

\(g\) adjoint à \deleted{\unsure{droite}}\unsure{gauche} de \(f\), donc on a

\diagram{\begin{tikzcd}[ampersand replacement=\&,column sep=0pt]
|[draw]| \operatorname{id}_A\xrightarrow{\alpha}gf \&
|[draw]| \underset{\substack{\Vert\\\varphi}}{fg}\xrightarrow{\beta}\operatorname{id}_B
\end{tikzcd}}

On a aussi

\diagram{\begin{tikzcd}[ampersand replacement=\&,row sep=small,column sep=small]
\&\&\lambda=f*\alpha*g\\
\boxed{\varphi\xrightarrow{\lambda}\varphi^2}\&\text{i.e.}\&
\deleted{\unsure{gf}}fg\longrightarrow \deleted{\unsure{g}}f\underbrace{gf}g\\
\&\&\operatorname{id}_A\arrow[u,"{\alpha}"']
\end{tikzcd}}

en plus de \(\boxed{\varphi\xrightarrow{\beta}\operatorname{id}_B}\).

\(\lambda\) et \(\beta\) permettent de définir, pour tout objet \(X\) de \(B\), une structure \(1/2\) \deleted{\unsure{covariant|cosimpliciale}} co-semisimpliciale \(\Phi^*(X)\)

\diagram{\begin{tikzcd}[ampersand replacement=\&,column sep=large,row sep=tiny]
\varphi(X)\arrow[r,shift left,"{\deleted{\unsure{\lambda}}\varphi*\beta}"]\arrow[r,shift right,"{\beta*\varphi}"']\arrow[d,equal]\&
\varphi^2(X)\arrow[r,shift left=1.5]\arrow[r]\arrow[r,shift right=1.5]\arrow[d,equal]\&
\varphi^3(X)\arrow[d,equal]\\
\Phi^0(X)\&\Phi^1(X)\&\Phi^2(X)
\end{tikzcd}}

i.e. on a un foncteur co-\(1/2\)-simplicial

\diagram{\begin{tikzcd}[ampersand replacement=\&,column sep=large,row sep=tiny]
\deleted{\unsure{\Phi}}\varphi\arrow[r,shift left,"{\unsure{\varphi*\beta|\varphi\beta}}"]\arrow[r,shift right,"{\beta*\varphi}"']\arrow[d,equal]\&
\varphi^2\arrow[r,shift left=2,"{\varphi^2*\beta\unsure{*u}}"]\arrow[r,"{\varphi*\beta*\varphi}" description]\arrow[r,shift right=2,"{\beta*\varphi^2}"']\arrow[d,equal]\&
\varphi^3\arrow[d,equal]\&\ldots\\
\Phi^0\&\Phi^1\&\Phi^2\&
\end{tikzcd}}

\uline{NB}\quad Les opérations bords n'utilisent que \(\beta\)

\hspace{3em}\rule{4em}{0.4pt}\quad de dégén.\quad\rule{7em}{0.4pt}\(\lambda\)

Pour \unsure{tt|tout} \(Y\) dans \(A\), on a une coaugmentation

\[
f(\deleted{X}Y)\longrightarrow\Phi^*(f(\deleted{X}Y))
\]

qui est une \uline{équivalence d'homotopie}, i.e.

\[
f\xrightarrow{\approx}\Phi^*\circ f
\]

l'homotopie est donnée par

\[
f\xrightarrow{f*\alpha}{\unsure{\Phi|\varphi}}^0\circ f=\varphi f=fgf
\]

Pour tout \(X\) (sous-entendu: on veut exprimer une structure suppl., \deleted{indiquant que} \inserted{\unsure{exprimant si possible} une donnée d'un isom.} \(X\simeq f(Y)\)), on considère les hom. possibles d'objets co\(1/2\) simpliciaux

\[
X\longrightarrow\Phi^*(X)
\]

\diagram{\begin{tikzcd}[ampersand replacement=\&]X\arrow[r,"u"]\&\varphi(X)\arrow[r,shift left,"{\varphi(u)}"]\arrow[r,shift right,"{\lambda}"']\&\varphi^2(X)\end{tikzcd}}

i.e.

\diagram{\begin{tikzcd}[ampersand replacement=\&,row sep=small]
X\arrow[r,dash]\arrow[d,"u"']\&{}\\
\varphi(X)\arrow[r,shift left]\arrow[r,shift right]\&\varphi^2(X)
\end{tikzcd}}
\scanpage{73}
% archive sheet number: 72

\noindent\fbox{\parbox{.9\linewidth}{Foncteurs fidèles transportables comme foncteurs \hfill oubli de structure.}}

Soit \(\deleted{\unsure{f}}p:\unsure{\mathcal E|E}\unsure{\longrightarrow|-}\unsure{B|\mathcal B}\) un foncteur. Si \(B\) est un groupoïde, \inserted{on a \unsure{l'équiv.|dit que}:} (i) \(p\) est fibrant \deleted{\unsure{surjectif} (i bis) à \unsure{ch.} des \unsure{ftrs|\(f^*\)}} \deleted{(ii) \(p\) cofibré \unsure{et} \illegible{1} \unsure{(iii bis) existe les} \(f^*\)} \inserted{i.e. \(p^{\unsure{o|0}}\) fibrant} (iii) \(\forall\,x\xrightarrow{u}y\) de \(B\) \deleted{\(\unsure{\beta}\longrightarrow\unsure{\beta}\)} et \(\xi\in\operatorname{Ob}\mathcal E_x,\ \exists\,\xi\xrightarrow[\unsure{\sim}]{v}\eta\) tels que \(p(v)=u\). \deleted{\(\operatorname{Fl}(p):\operatorname{Fl}(\mathcal E)\longrightarrow\operatorname{Fl}(B)\) surjectif.} Alors les flèches cartésiennes \inserted{sont les flèches cocartésiennes} de \(\mathcal E\) et sont \deleted{\unsure{exactmt}} ses isomorphismes.

Soit \(p:\mathcal E\longrightarrow B\) un foncteur. On dit que \(p\) est transportable si \deleted{\unsure{pour tte fl.}} \(\mathcal E\times_B B^{\mathrm{is}}\longrightarrow B^{\mathrm{is}}\) est fibrant (ou ce qui revient au \(=\), cofibrant).

Donc \(p\) est \deleted{fibrant} transp. ssi \(p^{\unsure{o|0}}\) l'est. \deleted{Flèches de transport \(=\)} flèches cartésiennes de \(\mathcal E\times_B B^{\mathrm{is}}\) sont les flèches cocartésiennes et sont les flèches isom. de \(\mathcal E\).

Soit \(p\) un foncteur; Alors \(p\) est isom. à un foncteur transportable. / Soient \(p:\mathcal E\unsure{-|\longrightarrow}B\) et \(p':\mathcal E\longrightarrow B\) deux foncteurs transportables isomorphes. Alors il existe une auto-équivalence \(\Phi\) \unsure{tel|de} \(\mathcal E\), \inserted{\uline{isom. à} \(\operatorname{id}_{\mathcal E}\),} telle que \(\Phi\) soit une \(B\)-équivalence \((\mathcal E,p)\unsure{-|\longrightarrow}(\mathcal E,p')\). /

\unsure{Donner} un foncteur \(p\) définit une \deleted{foncteur} \deleted{transportable}, \inserted{Catégorie sur \(B\),} \deleted{\unsure{et une cat.}} \deleted{\unsure{de fonct. fibrés}} \inserted{\unsure{fibrant|fibrée} sur \(\mathcal E^{\mathrm{is}}\),} unique à \(B\)-\unsure{équiva-|isogr.ïdes}
\scanpage{74}
% archive sheet number: 73

\unsure{conservatif} \(=\) fibres des \unsure{groupoïdes|propriétés}

fidèle
\scanpage{75}
% archive sheet number: 74

\unsure{lence|leurs} près.

Soit \(p\) un foncteur transportable. Alors \(p\) est fidèle ssi les catégories fibres \(\unsure{\mathcal E|E}_b\) sont des catégories ordonnées [donc si \(p\) est fibrant, \(\mathcal E\) est une catégorie fibrée en cat \inserted{\unsure{pré}/}ordonnées].

Pour que \(p\) soit conservatif, il f. et s. que les catégories fibres soient des groupoïdes, \deleted{[ou encore que tout morphisme de \(\mathcal E\times_{\unsure{B|\mathcal B}} B^{\mathrm{is}}\) soit cartésien,} Donc pour que \(p\) soit fidèle et conservatif, il f. et s. que les \(\mathcal E_b\) soient \uline{discrètes}.

Soit \(p\) un foncteur. \deleted{\uline{fibrant} \unsure{Cela}} \inserted{Conditions} équivalentes

a) \(p\) fibrant, fidèle, conservatif

b) \(p\) fibrant : fibres discrètes

c) \deleted{\(p\) \unsure{fibrant et} fidèle,} \inserted{\(p\) est fidèle et} / \(\forall\,X\in\operatorname{Ob}\mathcal E,\quad\mathcal E_{/X}\unsure{-|\longrightarrow} B_{/p(X)}\) est une équivalence

\margin{NB. \unsure{la deux.|le deux.} cond. de c) signifie que \(p\) est fibré en groupoïdes}

d) \(\exists\,F\in\widehat B\), et une \(B\)-équivalence

\[
\mathcal E\xrightarrow{\approx} B_{/F}
\]

Les foncteurs ``oubli de structure'' sont les types des foncteurs \deleted{\unsure{transp.|transportables}} \uuline{fidèles} \deleted{\uline{\unsure{conservatifs}}} \inserted{\uline{transportables}}. Ils ont un autre trait: les catégories fibres sont non \unsure{seult} préordonnées, mais ordonnées, i.e. \inserted{des} fibres \uline{réduites}, ce qui s'exprime \inserted{\unsure{mieux}} en disant \deleted{\unsure{soit}} que l'on a \uline{unicité} de relèvement dans la déf. de la transporta-
\scanpage{76}
% archive sheet number: 75

bilité. Si \(p:\unsure{\mathcal E|E}\unsure{-|\longrightarrow}\unsure{B|\mathcal B}\) ne satisfait pas cette condition, une modification triviale de \(\mathcal E\) (\unsure{mod équivalence|une équivalence|aut-équivalence|avec équivalence}) la rend satisfaite.

Partant d'un foncteur fidèle transportable, on définit \deleted{sur \(B^{\mathrm{is}}\) une catégorie} un \uline{foncteur} sur \(B^{\mathrm{is}}\)

\[
\begin{aligned}
b&\unsure{\mapsto|\sim}\text{ Classes d'équiv \unsure{à|:} isom près de }\mathcal E_b\\
 &=\text{``Structures d'espèce }p\text{ sur }b\text{''}
\end{aligned}
\quad\left|\begin{gathered}\text{à valeurs dans les}\\\text{ensembles}\\\inserted{\text{ordonnés}}\\
(\text{\uline{str. \(\pm\) fine}}]\end{gathered}\right.
\]

(via ``\uline{transport de structure}'') de façon qu'un objet de \(\mathcal E\) se reconstitue à isom \inserted{unique} près comme objet de \(B\) ``muni d'une structure d'espèce \(p\)''. De façon précise, ceci permet de reconstituer la catégorie \(\mathcal E^{\unsure{\mathrm{is}|\Delta}}\) avec son foncteur \(\mathcal E^{\mathrm{is}}\xrightarrow{p^{\mathrm{is}}}B^{\mathrm{is}}\).

Pour pouvoir reconstituer \(\mathcal E\) lui-\(=\), il faudrait que \(\mathcal E\) soit fibrant (\unsure{disons|cloisonné|clésé}) sur \(B\), de façon qu'on pourra définir la notion de ``structure image inverse'' d'une structure d'espèce \(p\) sur \(c\in\operatorname{Ob} B\), par une flèche \(b\unsure{-|\longrightarrow}c\)'', de façon à avoir \inserted{\(=\)} un foncteur \(\Sigma:B^{\unsure{o|0|\circ}}\longrightarrow(\text{ens ordonnés})\) [et non \unsure{seult} \deleted{\unsure{sur}} \(B^{\mathrm{is}}\longrightarrow(\text{ens ordonnés})\)]. On reconstitue ainsi \(\mathcal E\) à \(B\)-équivalence près par la connaissance de \(\Sigma\), et \(=\) à \(B\)-isom près si \(\mathcal E\) est réduite.

Ainsi, \unsure{mod|une} \(B\)-équiv. (\(B\)-isom), les catégories \uline{fibrées} sur \(B\), à projections \uline{fidèles} \inserted{(resp. et réduites)}, s'interprètent en termes de foncteurs \inserted{\unsure{contr.}} \(B^o\longrightarrow(\unsure{\text{Ens ord}|\text{ens. ord.}})\), et pour l'exprimer le langage des structures est assez suggestif.
\scanpage{77}
% archive sheet number: 76

Morphismes et \unsure{corresp|comorp}\\
Liste Problèmes

5 \unsure{p}
\scanpage{78}
% archive sheet number: 77

\unsure{Questions|Question} sur les topos. (Pour Mlle \unsure{Zunigery|Zunigey|Lucigey} ?).

1) Définition et étude des ``immersions'' de topos. \unsure{ou|de|des} Caractérisation ? Par condition \(f_*\) pl fid i.e. \(f^*f_*\xrightarrow{\sim}\operatorname{id}\) (i.e. \(f^*\) passage au quotient? Dans quelle mesure la \deleted{\unsure{NB|N.B.}} connaissance de \(f:\unsure{\mathcal U|U}\longrightarrow\unsure{\mathcal V|V}\) détermine-t-elle la ``partie loc fermée'' lui donnant naissance ? Plus précisément, si

% slanted
\margin{Théorie des \(f_!\) et \(f^!\)}

% slanted
\margin{\unsure{faux}}

\margin{OK. pour imm. ouverte et imm. fermée.}

\diagram{\begin{tikzcd}[ampersand replacement=\&,row sep=tiny]
\mathcal U\arrow[r,<->,"f"]\arrow[dr,<->,"{f'}"']\&\mathcal V\\
\deleted{\unsure{\mathcal U'|U'}}\arrow[r,draw=none]\&\mathcal V'
\end{tikzcd}}

sont des hom d'imm. de topos de \(=\) \deleted{\unsure{source}} \inserted{but} qui sont ``équivalents à \deleted{\unsure{source}} \inserted{but} fixé'' [i.e. \(\exists\) équivalence

\diagram{\begin{tikzcd}[ampersand replacement=\&,row sep=small]
\mathcal V\arrow[rr,"g"]\arrow[dr,"{f^*}"']\&\&\mathcal V'\arrow[dl,"{f'^*}"]\\
\&\mathcal U\&
\end{tikzcd}}

rendant comm. à isom près le diagramme], alors \(g\) est-elle déterminée à isom unique près ? [Notion de ``monomorphisme essentiel de Topos'' ??].

Caractérisation des imm ouvertes [par \(f^*\) commutant aux \(\underleftarrow{\lim}\) quelc., en plus de \(f_*\) pl fid ?].

\margin{\uline{\unsure{non|voir}}}

Caractérisation des imm fermées [par la condition \(F\longrightarrow f_*f^*F\) un épim pour \unsure{tt|\(H\)} faisceau \uline{abélien} \(F\) -- attention, \unsure{contre|autre} en dehors de ce cas, prendre \(F=\varnothing\) \(\ldots\) -- en plus de \(f_*\) pl fid.] ?

% slanted
\margin{Application \unsure{aux} ``bonnes corresp.'' pour \unsure{commutation} \deleted{\unsure{leurs}} \(H^i\) aux \(\underrightarrow{\lim}\)}

\textcircled{2} Etude des \(f\) tels que \(f^*\) commute aux \(\underleftarrow{\lim}\) quelc. Alors \(f^*\) \deleted{est-il de la forme \(\unsure{g_*}\)} admet-il un adjoint à \uline{gauche} \deleted{\unsure{\(g\)}} \(f_!\) ? Dans quelle mesure le cas du foncteur \unsure{canonique} \(f^*:\deleted{\unsure{V}}\mathcal U\longrightarrow\mathcal U_{/S}\) est-il typique ??

% slanted
\margin{\uline{oui}}

\margin{Il ne l'est pas, cf pts essentiels. Regarder le cas des commutation aux \uline{Hom} relatifs. \(\ldots\)}
\scanpage{79}
% archive sheet number: 78

\textcircled{3} Étude du cas \(f^*\) pl fidèle. Étude du cas où \(f_*\) commute aux \(\underrightarrow{\lim}\) \unsure{filtrées|finies}. [et aussi, à \unsure{certaines} \(\underrightarrow{\lim}\) ou transforme épim en épim \(\ldots\)] \inserted{-- fam. ép. en fam. ép.}

4) Étude du cas \(f_*\) pl fidèle (Cf \(1^\circ\)))

5) Notion de morphisme \uline{dominant} de topos (\(f_*(\varnothing)=\varnothing\)). \deleted{\unsure{6)|Plus}} Lien avec le fait \(f^*\) \uline{conservatif}. \deleted{\unsure{6)}} (Peut-on dire \unsure{qque} chose sur le cas \(f_*\) conservatif?)

\deleted{6) Donner une caractérisation des foncteurs \unsure{exact|exacts} à gauche \(\unsure{\mathcal U|U}\unsure{-|\longrightarrow}\unsure{\mathcal V|V}\) permettant un recollement à la Cartier-Artin. Cf Deligne)}

% slanted
\margin{\deleted{OK.}}

\deleted{7) Les Morphismes de \uuline{\(\mathcal U\)}-topos \inserted{\(\mathcal U,\mathcal V\)} font une \(\mathcal U\)-cat (i.e. les Hom sont ``petits''). \deleted{et} L'ens des classes, à isom près, de tels morphismes est-il petit ?}

% slanted
\margin{\uuline{non}}

8) a) Théorie des structures : la \unsure{Artin|Notion} des \unsure{'}topos noethériens ?

\diagram{\begin{tikzcd}[ampersand replacement=\&,column sep=huge]
(\mathrm{Top})\arrow[r,"{\text{cat. des fctrs fibres}}","{\scriptscriptstyle >}" marking]\&
(\mathrm{Cat})\arrow[l,bend left=20,"{\widehat C\unsure{\leftarrow|\sim|\mathrel{\leftarrow\!\shortmid}}C}"]
\end{tikzcd}}

\margin{ou y-a-t-il lieu de considérer qu'on s'insère dans une catégorie plus spécifique que (Cat)?}

\inserted{b) Topos où tt ss-objet est \unsure{sommande} (donc où la th. des supports est triviale)}

\inserted{c) Topos engendrés par les ss-objet de l'objet final.}

9) Déterminer la \uline{catégorie} \uline{Hom} de Topos en topos annelés dans certains cas habituels. (On doit obtenir de nombreuses catégories discrètes).

10) Topos modulaires. Catégories fibrées resp. Champs essentiellement représentables \(\ldots\).

\diagram{\begin{tikzcd}[ampersand replacement=\&,column sep=small]
\deleted{\text{11)}}\&\deleted{(\mathrm{Esp.})\unsure{\mathrm{red}}}\arrow[r]\&
\deleted{(\mathrm{Esp.}\ \text{\unsure{sobres}})}\arrow[r]\&
\deleted{(\mathrm{Top})}\arrow[l,bend left=30,"{\deleted{\substack{\text{2-adjoint : espace}\\\text{des foncteurs}\\\text{fibres}}}}"']
\end{tikzcd}}
\scanpage{80}
% archive sheet number: 79

\deleted{deux adjoints de \(f\) existent et \unsure{soient|sont} isom. Il suffirait de l'avoir pour \deleted{des} catégories \deleted{\unsure{non vides}} \(C,C'\), telles que \(f,g\) ne soient pas des \deleted{\unsure{équivalences|foncteurs}}\inserted{pl. fid.} (on passe ensuite à \(\widehat C,\widehat C'\)).}

\deleted{Peut-on \(=\) prendre \(C=C', f=g\), i.e. peut on trouver une catégorie \(C\), et un endom \(f:C\unsure{-|\longrightarrow}C\) qui ne soit pas \deleted{une équivalence}\inserted{pl. fid.}, et qui soit isom. à son adjoint (à deux \unsure{bords|côtés})?}

OK. On soit \(C\) une catégorie ayant un objet nul \deleted{\unsure{0}} \(e\), mais tel que le foncteur constant \deleted{\unsure{\(C\longrightarrow\)}} \(C\longrightarrow C\) défini par \(f(x)=e\) ne soit pas pl fidèle i.e. \(C\) non une catégorie \unsure{discrétable|discrète|discrètisable}, donc \(C\) contient un objet \(x\unsure{\ne|=}e\), [Ex. minimal :

\diagram{\begin{tikzcd}[ampersand replacement=\&]e\arrow[r,shift right,"v"']\&a\arrow[l,shift right,"u"']\arrow[loop right,"{\deleted{\unsure{vu}}w}"]\end{tikzcd}\,:}

deux objets \(e,a\), trois flèches en plus des flèches identiques \(u,v,vu=w\), relations \deleted{\unsure{à prio.}}\inserted{de} composition sont déterminées \unsure{uniqt|uniquement} alors par ces conditions, \deleted{sauf} une [\(uv=\operatorname{id}_e\quad uw=u,wv=v\)], savoir que il faut dire ce qu'est \(w^2\) : On peut prendre \(w^2=\operatorname{id}_a\) ou \(w^2=w\) au choix], On trouve le
